
% ==================== front ====================
% Самостоятельный конспект лекции 7.
% Сборка из корня курса: latexmk 7/rendezvous.tex
% Для оглавления и перекрёстных ссылок требуется повторный запуск pdflatex.
\documentclass[a4paper,11pt]{article}
\usepackage{iftex}
\RequirePDFTeX
\usepackage{cmap}
\usepackage[T2A]{fontenc}
\usepackage[utf8]{inputenc}
\usepackage[english,russian]{babel}
\usepackage{PTSerif}
\usepackage{amsmath,amssymb,amsthm,mathtools}
\usepackage[a4paper,left=20mm,right=20mm,top=22mm,bottom=23mm,headheight=14pt]{geometry}
\usepackage{microtype}
\usepackage{graphicx,xcolor,tikz,booktabs,tabularx,array,longtable}
\usetikzlibrary{arrows.meta,calc}
\usepackage{enumitem,fancyhdr,needspace,float,placeins}
\usepackage[most]{tcolorbox}
\newtcolorbox{explanation}[1]{enhanced,
 colback=accent!3!white,colframe=accent!80!black,
 boxrule=.5pt,toprule=1.2pt,arc=1mm,
 left=3mm,right=3mm,top=3mm,bottom=2mm,
 before skip=9pt,after skip=9pt,
 title={#1},fonttitle=\bfseries,coltitle=accent!70!black,
 attach boxed title to top left={xshift=3mm,yshift=-2mm},
 boxed title style={colback=white,frame hidden,boxrule=0pt,
 left=1.5mm,right=1.5mm,top=.4mm,bottom=.4mm},
 fontupper=\normalsize}
\newtcolorbox{remarkbox}{enhanced,
 colback=orange!4!white,colframe=orange!65!black,
 boxrule=.5pt,toprule=1.2pt,arc=1mm,
 left=3mm,right=3mm,top=2mm,bottom=2mm,
 before skip=9pt,after skip=9pt,fontupper=\normalsize}
\setcounter{topnumber}{2}
\renewcommand{\topfraction}{.85}
\renewcommand{\textfraction}{.12}
\renewcommand{\floatpagefraction}{.85}
\usepackage{hyperref}
\usepackage{bookmark}
\definecolor{accent}{HTML}{187AA5}
\definecolor{muted}{HTML}{5A6166}
\definecolor{figink}{HTML}{25323F}
\definecolor{figgray}{HTML}{687582}
\definecolor{figrad}{HTML}{176AA6}
\definecolor{figtan}{HTML}{087F78}
\definecolor{fignormal}{HTML}{B57714}
\definecolor{figplane}{HTML}{EEF4F8}
\definecolor{figorbit}{HTML}{79539C}
\hypersetup{unicode=true,colorlinks=true,linkcolor=accent,citecolor=accent,urlcolor=accent,
 pdftitle={Задача встречи и относительное движение},
 pdfauthor={Материалы курса «Введение в маневрирование КА»},
 pdfsubject={Общая и линейная встреча, задача Ламберта и модели относительного движения}}
\newcommand{\dd}{\mathop{}\!\mathrm d}
\newcommand{\R}{\mathbb R}
\DeclareMathOperator{\sgn}{sgn}
\numberwithin{equation}{section}
\newtheorem{theorem}{Теорема}[section]
\newtheorem{lemma}[theorem]{Лемма}
\newtheorem{proposition}[theorem]{Утверждение}
\theoremstyle{remark}
\newtheorem{remark}[theorem]{Замечание}
\renewcommand{\proofname}{Доказательство}
\setcounter{tocdepth}{3}
\setcounter{secnumdepth}{3}
\setlength{\parindent}{1.3em}
\setlength{\parskip}{2pt}
\linespread{1.055}
\setlist{itemsep=2pt,topsep=4pt,parsep=0pt}
\setlength{\emergencystretch}{1.5em}
\clubpenalty=10000
\widowpenalty=10000
\displaywidowpenalty=10000
\allowdisplaybreaks[1]
\pagestyle{fancy}
\fancyhf{}
\fancyhead[L]{\small\color{muted}Задача встречи и относительное движение}
\fancyfoot[C]{\small\thepage}
\renewcommand{\headrulewidth}{.3pt}
\renewcommand{\footrulewidth}{0pt}
\newcommand{\vect}[1]{\boldsymbol{#1}}
\newcommand{\norm}[1]{\left\lVert#1\right\rVert}
\newcommand{\mat}[1]{\boldsymbol{#1}}
\DeclareMathOperator{\atanTwo}{atan2}
\DeclareMathOperator{\rank}{rank}
\DeclareMathOperator{\diag}{diag}
\newcommand{\coursefigure}[3][.94]{%
\begin{figure}[htbp]\centering
\includegraphics[width=#1\textwidth]{images/#2.pdf}
\caption{#3}\label{fig:#2}\end{figure}}
% Generated by numerics.py. Values use km, s and classical elements.
\newcommand{\ExampleDuration}{1200\,\text{с}}
\newcommand{\ExampleInitialDistance}{2.859\,\text{км}}
\newcommand{\ExampleLinearCost}{4.4086\,\text{м/с}}
\newcommand{\ExampleLambertCost}{4.4086\,\text{м/с}}
\newcommand{\ExampleLinearMiss}{0.492\,\text{м}}
\newcommand{\ExampleLambertMiss}{<10^{-6}\,\text{м}}
\newcommand{\ExampleCorrectedMiss}{<10^{-6}\,\text{м}}
\newcommand{\ExampleLambertP}{7001.186\,\text{км}}
\newcommand{\ExampleTransferAngle}{74.437^\circ}
% Columns: method, |dV1| m/s, |dV2| m/s, sum m/s, position miss m.
\newcommand{\ExampleControlRows}{%
Линейное & 1.9393 & 2.4693 & 4.4086 & 0.4922 \\
Ламберт & 1.9394 & 2.4692 & 4.4086 & $<10^{-6}$ \\
Уточнённое & 1.9394 & 2.4692 & 4.4086 & $<10^{-6}$ \\
}
% Initial target elements, followed by chaser-minus-target deviations.
\newcommand{\ExampleInitialData}{%
a_0 &= 7000\,\text{км}, & e_0 &= 0.01, & i_0 &=45^\circ,\\
\Omega_0 &=20^\circ, & \omega_0 &=30^\circ, & M_0 &=40^\circ,\\
\delta a_0 &=-0.3\,\text{км}, & \delta e_0 &=2\cdot10^{-5}, & \delta i_0 &=10^{-4}\,\text{рад},\\
\delta\Omega_0 &=-1.5\cdot10^{-4}\,\text{рад}, & \delta\omega_0 &=10^{-4}\,\text{рад}, & \delta M_0 &=-4\cdot10^{-4}\,\text{рад}.
}
% Columns: method, sum dV (m/s), position miss (m), velocity miss (m/s).
\newcommand{\ExampleComparisonRows}{%
Линейное & 4.40864 & 0.4922 & $7.4\cdot10^{-4}$ \\
Ламберт & 4.40860 & $<10^{-6}$ & $<10^{-9}$ \\
Уточнённое & 4.40860 & $<10^{-6}$ & $<10^{-9}$ \\
}
% Columns: iteration, position miss (m), velocity miss (m/s), sum dV (m/s).
\newcommand{\ExampleCorrectionRows}{%
0 & 0.4922 & $7.4\cdot10^{-4}$ & 4.40864 \\
1 & $1.11\cdot10^{-6}$ & $<10^{-9}$ & 4.40860 \\
2 & $<10^{-6}$ & $<10^{-9}$ & 4.40860 \\
}
\newcommand{\ExampleLinearComponents}{%
$\Delta V_{1}$ & -1.5138 & 1.1890 & 0.2361 \\
$\Delta V_{2}$ & -2.0133 & -1.0124 & 1.0095 \\
}
\newcommand{\ExampleLambertComponents}{%
$\Delta V_{1}$ & -1.5140 & 1.1888 & 0.2363 \\
$\Delta V_{2}$ & -2.0133 & -1.0123 & 1.0094 \\
}
\newcommand{\ExampleCorrectedComponents}{%
$\Delta V_{1}$ & -1.5140 & 1.1888 & 0.2363 \\
$\Delta V_{2}$ & -2.0133 & -1.0123 & 1.0094 \\
}

\begin{document}
\begin{titlepage}
\centering
{\large Московский физико-технический институт\par}
\vspace{8mm}
{\large Курс «Введение в маневрирование КА»\par}
\vspace{32mm}
{\color{accent}\rule{28mm}{.8pt}\par}
\vspace{9mm}
{\LARGE\bfseries Задача встречи\\[3mm]
и относительное движение\par}
\vspace{10mm}
{\large Методические материалы\par}
\vspace{8mm}
{Общая и линейная постановки. Задача Ламберта.\\
Модели относительного движения и уточнение решения.\par}
\vfill
{Рабочая редакция\par}
\vspace{3mm}
{3 ноября 2026 года\par}
\end {titlepage}
\setcounter{page}{2}
\tableofcontents
\clearpage

\section*{Предмет рассмотрения и обозначения}
\addcontentsline{toc}{section}{Предмет рассмотрения и обозначения}

Задача встречи задаёт полное состояние аппарата в момент $t_f$,
совпадающее с состоянием цели. Рассматриваются общая задача
управления, её линейное импульсное приближение, задача Ламберта
и модели относительного движения CW, TH и SS.

Общие уравнения записываются в классических оскулирующих элементах
\[
 \vect y=(a,e,\omega,i,\Omega,M)^\top.
\]
В относительной системе используются декартовы координаты
и их производные во вращающемся базисе.
Импульс означает мгновенный скачок инерциальной скорости
при непрерывном положении. Его влияние на элементы
линеаризуется; проверочный прогноз выполняется
с точным скачком скорости.

\begin{center}
\renewcommand{\arraystretch}{1.13}
\begin{tabularx}{\textwidth}{@{}>{$}l<{$}X@{}}
\toprule
\multicolumn{1}{@{}l}{Обозначение}&Смысл\\\midrule
\mu & Гравитационный параметр центрального тела.\\
\vect R,\vect V & Инерциальные радиус-вектор и скорость.\\
p,h,n & Фокальный параметр $p=a(1-e^2)$, удельный момент
$h=\sqrt{\mu p}$, среднее движение $n=\sqrt{\mu/a^3}$.\\
\nu,u & Истинная аномалия и аргумент широты $u=\omega+\nu$.\\
S,T,W & Радиальная, трансверсальная и нормальная компоненты
управляющего ускорения в ОСК.\\
B,\Phi,M_j & Мгновенная чувствительность элементов к ускорению;
матрица изохронных производных; матрица манёвра.\\
\Delta\vect V & Столбец всех трёхкомпонентных импульсов.\\
\vect\rho=(x,y,z)^\top & Относительное положение в RTN цели:
радиально наружу, поперёк радиуса, по нормали к плоскости.\\
\varpi & Угол ориентации вектора эксцентриситета в плоскости
перелёта в разделе о задаче Ламберта.\\
\end{tabularx}
\end{center}

Углы в формулах измеряются в радианах. В таблицах исходных данных
их значения могут приводиться в градусах. Для импульсов на границах
интервала начальное состояние понимается до первого импульса,
а конечное --- после последнего. Все граничные сравнения изохронны:
состояния сопоставляются при одном физическом времени $t_f$.

\begin{remarkbox}
Классические элементы являются локальными координатами.
При $e=0$ не определены по отдельности $\omega$ и $M$,
а при $\sin i=0$ --- линия узлов. Поэтому общая матрица $B$
ниже используется при $e>0$, $\sin i\ne0$.
Малый ненулевой эксцентриситет не заменяется нулём в членах $1/e$.
Круговая опорная орбита в декартовых уравнениях относительного движения
такой координатной особенности не имеет.
\end{remarkbox}

\clearpage
\section{Общая постановка встречи в классических элементах}
\label{sec:general}

\subsection{Инерциальная динамика и орбитальная система координат}

Пусть $\vect a=(S,T,W)^\top$ --- компоненты управляющего
ускорения в ОСК. Уравнения управляемого аппарата в инерциальных осях
имеют вид
\begin{equation}
 \dot{\vect R}=\vect V,\qquad
 \dot{\vect V}=-\frac{\mu}{R^3}\vect R+Q\vect a,
 \qquad R=\norm{\vect R}.
 \label{eq:inertial}
\end{equation}
Единичные векторы орбитальной системы координат задаются как
\begin{equation}
 \vect e_r=\frac{\vect R}{R},\qquad
 \vect e_n=\frac{\vect R\times\vect V}{h},\qquad
 \vect e_t=\vect e_n\times\vect e_r,
 \quad h=\norm{\vect R\times\vect V}.
 \label{eq:RTN}
\end{equation}
Положим $Q=[\vect e_r\ \vect e_t\ \vect e_n]$. Компоненты ускорения в ОСК:
\begin{equation}
 \vect a=(S,T,W)^\top,\qquad Q\vect a=S\vect e_r+T\vect e_t+W\vect e_n.
 \label{eq:control-components}
\end{equation}
Вектор $\vect e_t$ поперечен радиусу в текущей плоскости орбиты.
На эллипсе он не совпадает с направлением скорости:
\[
 \vect V=V_r\vect e_r+V_t\vect e_t,
 \qquad
 V_r=\frac{\mu}{h}e\sin\nu,\qquad V_t=\frac{h}{R}.
\]

Классические элементы вычисляются по текущим $\vect R,\vect V$
и задают соответствующую им оскулирующую орбиту.
При кеплеровом движении первые пять элементов постоянны,
а $\dot M=n(a)$.

\subsection{Большая полуось}

Удельная механическая энергия равна
\begin{equation}
 \mathcal E=\frac{V^2}{2}-\frac{\mu}{R}=-\frac{\mu}{2a}.
 \label{eq:energy}
\end{equation}
Дифференцирование с учётом \eqref{eq:inertial} даёт
\[
 \dot{\mathcal E}=\vect V\cdot Q\vect a=V_rS+V_tT,
 \qquad
 \frac{\mu}{2a^2}\dot a=\frac{\mu}{h}e\sin\nu\,S+\frac{h}{R}T.
\]
Используя $h^2=\mu p$, получаем первую строку уравнений Гаусса:
\begin{equation}
 \dot a=\frac{2a^2}{h}
       \left(e\sin\nu\,S+\frac pR T\right).
 \label{eq:gauss-a}
\end{equation}
Нормальное ускорение ортогонально скорости и не входит в
\eqref{eq:gauss-a}.

\subsection{Эксцентриситет и направление перицентра}

Вектор удельного момента и вектор эксцентриситета определяются
соотношениями
\begin{equation}
 \vect h=\vect R\times\vect V,\qquad
 \vect e=\frac{\vect V\times\vect h}{\mu}-\vect e_r.
 \label{eq:invariants}
\end{equation}
Управляющее ускорение даёт
\begin{equation}
 \dot{\vect h}=\vect R\times Q\vect a,
 \qquad
 \dot{\vect e}=\frac{1}{\mu}
 \left((Q\vect a)\times\vect h+
              \vect V\times(\vect R\times Q\vect a)\right).
 \label{eq:edot-vector}
\end{equation}
Центральные гравитационные члены во второй формуле сокращаются.

Для получения скалярного $\dot e$ можно также использовать
$p=h^2/\mu=a(1-e^2)$. Из первой формулы
\eqref{eq:edot-vector} имеем $\dot h=RT$, поэтому
\[
 \dot p=\frac{2hR}{\mu}T
       =(1-e^2)\dot a-2ae\dot e.
\]
Подстановка \eqref{eq:gauss-a}, использование
$R=p/(1+e\cos\nu)$ и приведение членов дают
\begin{equation}
 \dot e=\frac1h\left[
       p\sin\nu\,S+
       \bigl((p+R)\cos\nu+Re\bigr)T\right].
 \label{eq:gauss-e}
\end{equation}
Проекция \eqref{eq:edot-vector} на направление, перпендикулярное
$\vect e$ в плоскости орбиты, даёт
\begin{equation}
 \dot\omega+\cos i\,\dot\Omega
 =\frac{1}{he}\left[-p\cos\nu\,S+(p+R)\sin\nu\,T\right].
 \label{eq:gauss-peri-inplane}
\end{equation}
Слагаемое $\cos i\,\dot\Omega$ учитывает поворот линии узлов.

\subsection{Поворот плоскости: наклонение и узел}

Из $\dot{\vect h}=\vect R\times Q\vect a$
следует поворот плоскости с угловой скоростью, пропорциональной
$RW/h$. Для $u=\omega+\nu$ получаем
\begin{equation}
 \dot i=\frac{R\cos u}{h}W,
 \qquad
 \dot\Omega=\frac{R\sin u}{h\sin i}W,
 \qquad u=\omega+\nu.
 \label{eq:gauss-plane}
\end{equation}
Из \eqref{eq:gauss-peri-inplane} и \eqref{eq:gauss-plane}
следует
\begin{equation}
 \dot\omega=\frac{-p\cos\nu\,S+(p+R)\sin\nu\,T}{he}
                  -\frac{R\sin u\cot i}{h}W.
 \label{eq:gauss-omega}
\end{equation}
Члены $1/e$ и $1/\sin i$ связаны с особенностями классических
элементов при круговой и экваториальной орбитах.

\subsection{Изменение средней аномалии}

Используется текущая средняя аномалия $M$.
Между истинной аномалией $\nu$,
эксцентрической аномалией $E$ и $M$ имеются связи
\begin{equation}
 M=E-e\sin E,\qquad
 \tan\frac E2=\sqrt{\frac{1-e}{1+e}}\tan\frac\nu2.
 \label{eq:anomalies}
\end{equation}
Для сохранения квадранта $E$ вычисляют двухаргументным арктангенсом.
Производные $M$ при независимых $e,\nu$ равны
\begin{equation}
 \left.\frac{\partial M}{\partial\nu}\right|_e
 =\frac{(1-e^2)^{3/2}}{(1+e\cos\nu)^2},\qquad
 \left.\frac{\partial M}{\partial e}\right|_\nu
 =-\frac{\sqrt{1-e^2}\sin\nu(2+e\cos\nu)}
              {(1+e\cos\nu)^2}.
 \label{eq:M-derivatives}
\end{equation}
В оскулирующем описании
$\dot\nu=h/R^2-\dot\omega-\cos i\dot\Omega$.
Дифференцируем $M(e,\nu)$, подставляем
\eqref{eq:gauss-e}, \eqref{eq:gauss-peri-inplane} и
используем $h/R^2$ для кеплеровой части. Получаем
\begin{equation}
 \dot M=n-\frac{2R}{na^2}S
       -\sqrt{1-e^2}(\dot\omega+\cos i\dot\Omega).
 \label{eq:gauss-M-compact}
\end{equation}
В выражении для $\dot M$ коэффициент при нормальном ускорении
$W$ равен нулю. В развёрнутом виде
\begin{equation}
 \dot M=n+\frac{\sqrt{1-e^2}}{he}
 \left[(p\cos\nu-2Re)S-(p+R)\sin\nu\,T\right].
 \label{eq:gauss-M}
\end{equation}
Зависимость $n$ от $a$ учитывается далее в матрице перехода.

\subsection{Матрица чувствительности \texorpdfstring{$B$}{B}}

Соберём полученные формулы в матричную запись
\begin{equation}
 \dot{\vect y}=\vect f_0(\vect y)+B(\vect y)\vect a,\qquad
 \vect f_0=(0,0,0,0,0,n(a))^\top,\qquad
 \vect a=(S,T,W)^\top.
 \label{eq:gauss-matrix}
\end{equation}
Для порядка $(a,e,\omega,i,\Omega,M)$ имеем
\begin{equation}
 B(\vect y)=\begin{pmatrix}
 \dfrac{2a^2e\sin\nu}{\sqrt{\mu p}}&
 \dfrac{2a^2(1+e\cos\nu)}{\sqrt{\mu p}}&0\\[5pt]
 \sqrt{\dfrac p\mu}\sin\nu&
 \sqrt{\dfrac p\mu}
 \dfrac{2\cos\nu+e(1+\cos^2\nu)}{1+e\cos\nu}&0\\[5pt]
 -\sqrt{\dfrac p\mu}\dfrac{\cos\nu}{e}&
 \sqrt{\dfrac p\mu}
 \dfrac{(2+e\cos\nu)\sin\nu}{e(1+e\cos\nu)}&
 -\sqrt{\dfrac p\mu}\dfrac{\sin u\cot i}{1+e\cos\nu}\\[5pt]
 0&0&\sqrt{\dfrac p\mu}\dfrac{\cos u}{1+e\cos\nu}\\[5pt]
 0&0&\sqrt{\dfrac p\mu}\dfrac{\sin u}{(1+e\cos\nu)\sin i}\\[5pt]
 \dfrac{p}{e\sqrt{\mu a}}
 \left(\cos\nu-\dfrac{2e}{1+e\cos\nu}\right)&
 -\dfrac{p(2+e\cos\nu)\sin\nu}{e\sqrt{\mu a}(1+e\cos\nu)}&0
 \end{pmatrix}.
 \label{eq:B-classical}
\end{equation}
Здесь $p=a(1-e^2)$, $u=\omega+\nu$.
Истинная аномалия $\nu$ восстанавливается из $M$ уравнением
Кеплера. Матрица $B$ меняется вдоль орбиты;
её столбцы соответствуют компонентам $S,T,W$.

При учёте возмущений уравнения имеют вид
\begin{equation}
 \dot{\vect y}=\vect f_0(t,\vect y)+B(\vect y)\vect a.
 \label{eq:general-element-dynamics}
\end{equation}
Возмущающие ускорения учитываются в $\vect f_0$;
матрица $B$ сохраняется. Постоянные компоненты $(S,T,W)^\top$
во вращающейся ОСК соответствуют переменному инерциальному
вектору $Q\vect a$.

\subsection{Краевая задача и критерий выбора управления}

Пусть управляемый аппарат имеет начальное состояние $\vect y_i$
в момент $t_i$, а цель --- известное движение
$\vect y_{\rm tar}(t)$. Если цель движется без управления, её
конечное состояние получается отдельным прогнозом из её
начальных элементов. Если задана эфемерида, берётся значение
из неё. В обоих случаях в конечном условии стоит
$\vect y_t=\vect y_{\rm tar}(t_f)$, относящееся именно к $t_f$.

Общая задача встречи формулируется следующим образом:
\begin{equation}
 \begin{aligned}
  &\min_{\vect a\in\mathcal A}
     J=\int_{t_i}^{t_f}\norm{\vect a(t)}\,\dd t,\\
  &\dot{\vect y}=\vect f_0(t,\vect y)+B(\vect y)\vect a(t),\\
  &\vect y(t_i)=\vect y_i,\qquad \vect y(t_f)=\vect y_t.
 \end{aligned}
 \label{eq:general-rendezvous}
\end{equation}
Множество $\mathcal A$ задаёт допустимое управление: ограничения
модуля ускорения, интервалы разрешённой работы, ориентационные
и другие ограничения конкретной задачи. В \eqref{eq:general-rendezvous}
времена $t_i,t_f$ фиксированы. Если какое-либо время выбирается,
оно должно быть явно включено в список переменных.

Интеграл нормы ускорения называется характеристической скоростью.
Полное конечное условие содержит фазу $M$;
равенство только первых пяти элементов задаёт геометрическую орбиту.
Общая постановка следует \cite{chiniaev} и здесь записана
в классических элементах.

\FloatBarrier
\section{Линейная импульсная задача встречи}
\label{sec:linear}

\subsection{Импульс как локальное изменение элементов}

В импульсной модели положение непрерывно, а инерциальная скорость
испытывает скачок:
\begin{equation}
 \vect R_j^+=\vect R_j^-,\qquad
 \vect V_j^+=\vect V_j^-+Q_j\Delta\vect V_j,
 \qquad Q_j=[\vect e_{rj}\ \vect e_{tj}\ \vect e_{nj}].
 \label{eq:exact-impulse}
\end{equation}
Здесь $\Delta\vect V_j$ --- компоненты импульса в предманёвренной
ОСК. Формула \eqref{eq:exact-impulse} является точным определением
идеального импульса. Орбитальные элементы после него находятся
по новому положению и скорости.

Однако зависимость элементов от скорости нелинейна. Например,
из энергии непосредственно следует
\[
 \mathcal E^+-\mathcal E^-
 =\vect V^-\cdot Q_j\Delta\vect V_j
             +\frac12\norm{\Delta\vect V_j}^2.
\]
Последний член уже не описывается первой производной.
При малом импульсе из уравнений Гаусса получается
\begin{equation}
 \vect y_j^+-\vect y_j^-
 =B(\vect y_j^-)\Delta\vect V_j
                  +O(\norm{\Delta\vect V_j}^2).
 \label{eq:linear-impulse}
\end{equation}

\subsection{Движение без управления и конечное отклонение}

Из заданного начального состояния вычислим траекторию без
управления $\vect y_0(t)$:
\begin{equation}
 \dot{\vect y}_0=\vect f_0(t,\vect y_0),\qquad
 \vect y_0(t_i)=\vect y_i.
 \label{eq:reference}
\end{equation}
Потребное изменение конечного
состояния равно
\begin{equation}
 \Delta\vect y_f=\vect y_t-\vect y_0(t_f).
 \label{eq:terminal-deficit}
\end{equation}
В \eqref{eq:terminal-deficit} сравниваются состояния при одном $t_f$.
Чтобы получить линейное приближение, положим
$\vect y=\vect y_0+\delta\vect y$ и оставим первый порядок
по отклонениям и управлению:
\begin{equation}
 \delta\dot{\vect y}=A(t)\delta\vect y+B_0(t)\vect a,
 \quad
 A(t)=\left.\frac{\partial\vect f_0}{\partial\vect y}\right|_{\vect y_0(t)},
 \quad B_0(t)=B(\vect y_0(t)).
 \label{eq:variation}
\end{equation}
Произведение изменения $B$ на управление имеет второй порядок.
Все коэффициенты вычисляются на одной опорной траектории.

\subsection{Матрица изохронных производных}

Матрица перехода $\Phi(t,t_j)$ определяется уравнением
\begin{equation}
 \frac{\dd\Phi(t,t_j)}{\dd t}=A(t)\Phi(t,t_j),
 \qquad \Phi(t_j,t_j)=I_6.
 \label{eq:stm-definition}
\end{equation}
Матрица переносит отклонение между заданными моментами:
$\delta\vect y(t)=\Phi(t,t_j)\delta\vect y(t_j)$.
Она удовлетворяет $\Phi(t_3,t_1)=\Phi(t_3,t_2)\Phi(t_2,t_1)$.

В кеплеровом случае закон движения известен точно:
\[
 \vect y(t)=
 (a_j,e_j,\omega_j,i_j,\Omega_j,M_j+n(a_j)(t-t_j))^\top.
\]
Дифференцирование по состоянию в $t_j$ даёт
\begin{equation}
 \Phi(t,t_j)=
 \begin{pmatrix}
 1&0&0&0&0&0\\0&1&0&0&0&0\\0&0&1&0&0&0\\
 0&0&0&1&0&0\\0&0&0&0&1&0\\
 -\dfrac{3n}{2a}(t-t_j)&0&0&0&0&1
 \end{pmatrix},\qquad
 \frac{\dd n}{\dd a}=-\frac{3n}{2a}.
 \label{eq:kepler-stm}
\end{equation}
Приращение большой полуоси меняет среднее движение,
поэтому конечная фаза зависит от времени переноса:
\begin{equation}
 \delta M(t)=\delta M_j-\frac{3n}{2a}(t-t_j)\delta a_j.
 \label{eq:phase-variation}
\end{equation}
\coursefigure{phi-phase}{Изменение фазы при $a=7000$~км,
$\delta a=100$~км на интервале 24~ч. Сравниваются состояния
в одинаковые моменты времени; наклон линейной зависимости
задаётся $\partial n/\partial a=-3n/(2a)$.}

Матрица \eqref{eq:kepler-stm} --- точный якобиан кеплерова
отображения. Для конечного приращения остаётся член второго порядка:
\[
 n(a+\delta a)-n(a)
 =-\frac{3n}{2a}\delta a+\frac{15n}{8a^2}(\delta a)^2
 +O((\delta a)^3).
\]

\subsection{Матрица манёвра}

Совместим мгновенное изменение \eqref{eq:linear-impulse}
с переносом отклонения до $t_f$. Для одного импульса получаем
\begin{equation}
 \delta\vect y_{f,j}=M_j\Delta\vect V_j,
 \qquad
 M_j=\Phi(t_f,t_j)B(\vect y_0(t_j))\in\R^{6\times3}.
 \label{eq:maneuver-matrix}
\end{equation}
Матрица $B_j$ задаёт мгновенное изменение элементов,
$\Phi$ переносит его к $t_f$.

В кеплеровой модели первые пять строк $M_j$ совпадают
с первыми пятью строками $B_j=B(\vect y_0(t_j))$, а последняя
строка имеет особенно простой вид:
\begin{equation}
 (M_j)_{6,:}=(B_j)_{6,:}
       -\frac{3n}{2a}(t_f-t_j)(B_j)_{1,:}.
 \label{eq:last-row}
\end{equation}
Первая часть фазовой строки задаёт мгновенное изменение $M$,
вторая --- накопленное изменение среднего движения.

\subsection{Произвольное число импульсов и суперпозиция}

Пусть $N$ импульсов приложены в заранее заданные моменты
$t_i\le t_1<\ldots<t_N\le t_f$. В первом порядке вклады складываются:
\begin{equation}
 \sum_{j=1}^N M_j\Delta\vect V_j=\Delta\vect y_f.
 \label{eq:sum-boundary}
\end{equation}
Введём столбец управления и общую матрицу
\begin{equation}
 \Delta\vect V=
 \begin{pmatrix}\Delta\vect V_1\\\vdots\\\Delta\vect V_N\end{pmatrix}
 \in\R^{3N},\qquad
 \mathcal M=[M_1\ \cdots\ M_N]\in\R^{6\times3N}.
 \label{eq:stack-control}
\end{equation}
Тогда линейное конечное условие встречи записывается как
\begin{equation}
 \boxed{\mathcal M\Delta\vect V=\Delta\vect y_f.}
 \label{eq:linear-boundary}
\end{equation}
Это обобщение двухимпульсной записи \cite{chiniaev}.
Суммирование допустимо в первом порядке: изменения состояния
и базисов последующих манёвров дают произведения малых величин.

\subsection{Существование, единственность и обусловленность}

Равенство \eqref{eq:linear-boundary} разрешимо тогда и только тогда,
когда $\Delta\vect y_f$ принадлежит образу $\mathcal M$.
Для одного трёхкомпонентного импульса ранг не превышает трёх,
поэтому произвольный шестимерный вектор обычно недостижим.
При двух импульсах матрица квадратная; единственное решение
существует при её невырожденности. При $N>2$ возможно семейство
решений, но из одного лишь числа столбцов полный ранг не следует.

Близость столбцов к линейной зависимости повышает
чувствительность управления к исходным данным.
Для численного анализа строку $a$ и угловые строки необходимо
масштабировать. Иначе формальное число обусловленности зависит
от выбора километров или метров, радиан или градусов.
Пусть $D_y$ --- диагональная матрица характерных масштабов
состояния. Тогда анализируется
\[
 \widehat{\mathcal M}=D_y^{-1}\mathcal M D_V,
 \qquad
 \Delta\widehat{\vect y}_f=D_y^{-1}\Delta\vect y_f,
\]
где $D_V$ задаёт масштаб скорости.
При полном ранге масштабирование не меняет физическое равенство,
но делает численные оценки содержательными.

\coursefigure{rendezvous-functional-slide}{Характеристическая скорость линейной
двухимпульсной встречи в зависимости от развёрнутых истинных долгот
$L_1,L_2$, где $L=\Omega+\omega+\nu$.
Опорная орбита: $a=6800$~км, $e=0{,}01$, $i=45^\circ$,
$\omega=\Omega=M_i=0$; $t_f-t_i=24$~ч.
Конечное отклонение
$\Delta\vect y_f=(1\ \text{км},\ 0{,}01,\ 1^\circ,\ 0{,}1^\circ,\ 0{,}1^\circ,\ 1^\circ)^\top$
задано в порядке $(a,e,\omega,i,\Omega,M)$.
Белые точки --- найденные локальные минимумы; зелёные --- 50 лучших из найденных.}

\subsection{Допустимое решение и минимум характеристической скорости}

Для выбора управления при фиксированных моментах импульсов введём критерий:
\begin{equation}
 \begin{aligned}
  &\min_{\Delta\vect V}\ J=\sum_{j=1}^N\norm{\Delta\vect V_j},\\
  &\mathcal M\Delta\vect V=\Delta\vect y_f,
  \qquad \norm{\Delta\vect V_j}\le\Delta V_{\max,j}.
 \end{aligned}
 \label{eq:fixed-time-cost}
\end{equation}
Сумма евклидовых норм выпукла, а равенство линейно.
При дополнительных выпуклых ограничениях задача остаётся
выпуклой. Её можно записать как коническую задачу второго порядка,
введя $s_j\ge0$ и ограничения $\norm{\Delta\vect V_j}\le s_j$
с линейным функционалом $\sum_js_j$.

Другая распространённая задача минимизирует
$\sum_j\norm{\Delta\vect V_j}^2=\norm{\Delta\vect V}^2$.
Если $\mathcal M$ имеет полный строчный ранг, её решение равно
\begin{equation}
 \Delta\vect V_{\rm minnorm}
 =\mathcal M^\top(\mathcal M\mathcal M^\top)^{-1}
                    \Delta\vect y_f.
 \label{eq:minnorm}
\end{equation}
Для решения системы используют QR- или сингулярное разложение.
Минимум суммы квадратов и минимум суммы норм импульсов
в общем случае различаются. При двух фиксированных импульсах
и невырожденной матрице допустимое решение единственно;
выбор моментов импульсов остаётся отдельной задачей.

\subsection{Выбор моментов импульсов и границы линейности}

Если моменты манёвров становятся переменными, матрица сама
зависит от искомых величин:
\begin{equation}
 \mathcal M=\mathcal M(t_1,\ldots,t_N),\qquad
 M_j=\Phi(t_f,t_j)B(\vect y_0(t_j)).
 \label{eq:free-epochs}
\end{equation}
Моменты входят в $B(\vect y_0(t_j))$ и в длительность переноса.
Поэтому совместный выбор векторов импульсов и моментов
их приложения является нелинейной задачей. При выборе $t_f$ также пересчитывается
$\Delta\vect y_f(t_f)$. Если манёвры задаются событиями
по аномалии или долготе, их времена меняются вместе
с управлением и требуют соответствующих производных.
Полученное управление проверяется по исходным нелинейным уравнениям.

\FloatBarrier
\section{Задача Ламберта через геометрию семейства коник}
\label{sec:lambert}

\subsection{Постановка задачи}

В задаче Ламберта рассматривается один пассивный участок
в центральном поле. Заданы два положения $\vect R_1,\vect R_2$,
гравитационный параметр $\mu$ и положительная длительность
$\tau=t_2-t_1$. Требуется найти скорости $\vect V_1,\vect V_2$
на кеплеровой дуге, проходящей через эти положения за время $\tau$:
\begin{equation}
 \ddot{\vect R}=-\frac{\mu}{R^3}\vect R,
 \qquad \vect R(t_1)=\vect R_1,\qquad
 \vect R(t_2)=\vect R_2.
 \label{eq:lambert-boundary}
\end{equation}
Заданы также направление обхода и число полных оборотов $m$.
Для неколлинеарных радиус-векторов плоскость определяется
концами, но имеет два ориентированных нормальных направления.
Они задают обход с углом меньше или больше $\pi$.
Две многовитковые ветви при фиксированном $m$ рассматриваются отдельно.

В этом разделе параметр $p$ всегда означает фокальный параметр
коники. Линия изложения соответствует общей геометрической
структуре краевой задачи из \cite[\S~5.2]{bmw}; параметризация
через компоненты вектора эксцентриситета позволяет непосредственно
построить её однопараметрическое семейство.

\subsection{Семейство коник}

Выберем ортонормированный базис в плоскости перелёта:
\begin{equation}
 \vect e_1=\frac{\vect R_1}{R_1},\qquad
 \vect e_3=\pm\frac{\vect R_1\times\vect R_2}
                         {\norm{\vect R_1\times\vect R_2}},
 \qquad \vect e_2=\vect e_3\times\vect e_1.
 \label{eq:lambert-plane}
\end{equation}
Знак нормали выбирается в соответствии с направлением движения.
Тогда
\[
 \vect R_2=R_2(\cos\theta\,\vect e_1+\sin\theta\,\vect e_2),
 \qquad 0<\theta<2\pi.
\]
Здесь отдельно исключён случай $\sin\theta=0$:
для сонаправленных или противоположных радиусов стандартная
параметризация ниже вырождается.

Запишем вектор эксцентриситета в выбранном базисе:
\begin{equation}
 \vect e=e_c\vect e_1+e_s\vect e_2,\qquad
 e=\sqrt{e_c^2+e_s^2},\qquad
 \varpi=\atanTwo(e_s,e_c).
 \label{eq:lambert-evector}
\end{equation}
Полярное уравнение коники с фокусом в начале координат имеет вид
\begin{equation}
 R(\eta)=\frac{p}{D(\eta;p)},\qquad
 D(\eta;p)=1+e_c\cos\eta+e_s\sin\eta.
 \label{eq:conic-family}
\end{equation}
Угол $\eta$ отсчитывается от первого радиуса, поэтому
истинная аномалия данной коники равна $\nu=\eta-\varpi$.
Требования прохождения через два заданных положения дают
два скалярных равенства:
\begin{equation}
 1+e_c=\frac p{R_1},\qquad
 1+e_c\cos\theta+e_s\sin\theta=\frac p{R_2}.
 \label{eq:evector-constraints}
\end{equation}
Если выбрать $p$, остальные два неизвестных определяются явно:
\begin{equation}
 \boxed{e_c(p)=\frac p{R_1}-1,\qquad
 e_s(p)=\frac{p/R_2-1-e_c(p)\cos\theta}{\sin\theta}.}
 \label{eq:e-of-p}
\end{equation}
Следовательно, два положения оставляют один непрерывный
параметр семейства. Он определяет и форму, и ориентацию
соответствующей коники. \coursefigure{lambert-family}{Кеплеровы коники с общим фокусом
при $R_1=7000$~км, $R_2=10000$~км, $\theta=100^\circ$.
Выделено решение для времени перелёта 30~мин. Изменение
фокального параметра $p$ меняет вектор эксцентриситета,
ориентацию и время; точки на концах остаются прежними.}

По найденным $p,e$ определяется вид коники:
$e<1$ соответствует эллипсу, $e=1$ --- параболе,
$e>1$ --- гиперболе. При $e\ne1$ большая полуось равна
\begin{equation}
 a=\frac{p}{1-e^2}.
 \label{eq:a-from-p}
\end{equation}
Для эллипса $a>0$, для гиперболы $a<0$.
Вблизи параболы деление на $1-e^2$ плохо обусловлено;
само семейство \eqref{eq:conic-family} при этом не исчезает.

\subsection{Область допустимых значений}

Не всякое формальное значение $p$ задаёт нужный перелёт.
Необходимо $p>0$ и
\begin{equation}
 D(\eta;p)>0\quad
 \text{на всей дуге }0\le\eta\le\theta+2\pi m.
 \label{eq:admissible-arc}
\end{equation}
Условие проверяется на всей дуге. Нулевой знаменатель
соответствует асимптоте, отрицательный радиус не принадлежит
нужной физической ветви. При $e<1$ знаменатель всегда положителен;
при $e\ge1$ дуга ограничена асимптотами.
Полные обороты возможны только для эллипса.
Ограничение $R(\eta)>R_{\min}$ проверяется отдельно
по концам и принадлежащим дуге перицентрам.

\begin{explanation}{Ветви при заданной большой полуоси}
Можно описывать семейство через $a$, но при одном $a$
через заданные концы в общем случае проходят разные
эллиптические дуги. Их времена различны.
Поэтому выражение $T(a)$ без указания ветви неоднозначно.
Параметризация \eqref{eq:e-of-p} сохраняет конкретную
ориентированную конику, а число оборотов задаётся отдельно.
\end{explanation}

\subsection{Уравнение времени перелёта}

Сохранение удельного момента даёт
\[
 R^2\dot\eta=h=\sqrt{\mu p},\qquad
 \dd t=\frac{R^2}{\sqrt{\mu p}}\,\dd\eta.
\]
Поэтому время движения вдоль допустимой коники равно
\begin{equation}
 T_m(p)=\frac1{\sqrt{\mu p}}
  \int_0^{\theta+2\pi m}R(\eta;p)^2\,\dd\eta
 =\frac{p^{3/2}}{\sqrt\mu}
  \int_0^{\theta+2\pi m}\frac{\dd\eta}{D(\eta;p)^2}.
 \label{eq:lambert-time-integral}
\end{equation}
Каждая допустимая коника задаёт время движения по закону площадей.
Условие заданной длительности превращается в одно скалярное
уравнение
\begin{equation}
 F_m(p)=T_m(p)-\tau=0.
 \label{eq:lambert-scalar}
\end{equation}
Геометрические условия сокращают непрерывный поиск до одной
переменной $p$; направление и число оборотов остаются заданными.

\coursefigure{lambert-branches}{Время движения как функция
фокального параметра при $R_1=7000$~км, $R_2=10000$~км,
$\theta=100^\circ$. Горизонтальная линия задаёт 200~мин.
Минимум одновитковой ветви составляет около 156.8~мин:
выше него существуют два решения с одним полным оборотом.
Отсев дуг, пересекающих Землю, при построении не выполнен.}

\subsection{Вычисление времени перелёта}

Время можно вычислять квадратурой \eqref{eq:lambert-time-integral}
или аналитически через аномалии.
Положим $\nu_1=-\varpi$, $\nu_2=\theta-\varpi+2\pi m$.
На эллипсе
\begin{equation}
 \cos E=\frac{e+\cos\nu}{1+e\cos\nu},\qquad
 \sin E=\frac{\sqrt{1-e^2}\sin\nu}{1+e\cos\nu}.
 \label{eq:lambert-E}
\end{equation}
Значения $E$ непрерывно продолжаются вдоль дуги
с сохранением числа оборотов. Тогда
\begin{equation}
 T_m=\sqrt{\frac{a^3}{\mu}}
 \bigl[E_2-E_1-e(\sin E_2-\sin E_1)\bigr].
 \label{eq:lambert-time-ellipse}
\end{equation}
Сведение обеих аномалий по модулю $2\pi$ перед вычитанием
удаляет целые периоды и решает другую задачу.

Для гиперболы удобно ввести гиперболическую аномалию $H$:
\begin{equation}
 \sinh H=\frac{\sqrt{e^2-1}\sin\nu}{1+e\cos\nu},\qquad
 T=\sqrt{\frac{(-a)^3}{\mu}}
    \left[e(\sinh H_2-\sinh H_1)-(H_2-H_1)\right].
 \label{eq:lambert-time-hyperbola}
\end{equation}
Знаменатель должен быть положительным на всей дуге.
Для параболического случая используется формула Баркера.
При $D_B=\tan(\nu/2)$ она имеет вид
\begin{equation}
 T=\frac12\sqrt{\frac{p^3}{\mu}}
 \left[D_{B2}+\frac{D_{B2}^3}{3}
             -D_{B1}-\frac{D_{B1}^3}{3}\right].
 \label{eq:barker}
\end{equation}
Вблизи $e=1$ угловой интеграл позволяет избежать
вычитания больших близких величин в формулах через $a$.

\subsection{Поиск корня и учёт ветвей}

Для заданных $\theta,m$ выбирают допустимый интервал $p$,
вычисляют $e_c,e_s,e$ и время $T_m(p)$.
На каждой ветви находят интервал, на концах которого
$T_m(p)-\tau$ имеет разные знаки. Затем уточняют корень
методом бисекции или методом Ньютона с сохранением
приближений внутри интервала.

При необходимости производная времени получается
дифференцированием \eqref{eq:lambert-time-integral}:
\begin{equation}
 \frac{\dd T_m}{\dd p}
 =\frac{3T_m}{2p}
  -\frac{2p^{3/2}}{\sqrt\mu}
   \int_0^{\theta+2\pi m}
           \frac{D_p(\eta;p)}{D(\eta;p)^3}\,\dd\eta,
 \label{eq:time-derivative}
\end{equation}
где
\[
 D_p=\frac{\cos\eta}{R_1}
  +\frac{1/R_2-\cos\theta/R_1}{\sin\theta}\sin\eta.
\]
Производная определена внутри допустимого интервала.
Для невырожденной задачи без дополнительных оборотов
при фиксированном направлении обхода имеется одна
кеплерова дуга с заданным положительным временем.
С ограничениями на минимальный радиус или энергию она может
оказаться неприемлемой. Для заданного $m\ge1$ время имеет
минимум: при $\tau<T_{m,\min}$ решений нет, при равенстве
ветви сливаются, при большем времени имеются две ветви.
Обе ветви ищут в отдельных интервалах \cite{izzo}.

\subsection{Восстановление скоростей через \texorpdfstring{$f,g$}{f,g}}

Пусть найдено допустимое $p$, удовлетворяющее
\eqref{eq:lambert-scalar}. Связь положений и скоростей
на кеплеровой дуге имеет вид
\begin{equation}
 \vect R_2=f\vect R_1+g\vect V_1,\qquad
 \vect V_2=\dot f\vect R_1+\dot g\vect V_1.
 \label{eq:fg-relation}
\end{equation}
Для известных $p,R_1,R_2,\theta$ достаточно вычислить
\begin{equation}
 f=1-\frac{R_2}{p}(1-\cos\theta),\qquad
 g=\frac{R_1R_2\sin\theta}{\sqrt{\mu p}},\qquad
 \dot g=1-\frac{R_1}{p}(1-\cos\theta).
 \label{eq:fg-coefficients}
\end{equation}
Из первого равенства \eqref{eq:fg-relation} и тождества
$f\dot g-\dot f g=1$ получаем
\begin{equation}
 \boxed{\vect V_1=\frac{\vect R_2-f\vect R_1}{g},\qquad
 \vect V_2=\frac{\dot g\vect R_2-\vect R_1}{g}.}
 \label{eq:lambert-velocities}
\end{equation}
Скорости восстанавливаются алгебраически.
Есть и прямая геометрическая проверка этих скоростей.
В точке с углом $\eta$ радиальная и трансверсальная
компоненты равны
\begin{equation}
 V_r=\sqrt{\frac\mu p}
          (e_c\sin\eta-e_s\cos\eta),\qquad
 V_t=\frac{\sqrt{\mu p}}{R(\eta)}.
 \label{eq:conic-velocities}
\end{equation}
При $\theta\to0$ или $\pi$ коэффициент $g$ мал.
Строго коллинеарные концы требуют дополнительного выбора
плоскости и предельных формул.

\subsection{Включение дуги Ламберта в задачу встречи}

При граничных импульсах в $t_i,t_f$ первая точка дуги равна
начальному положению управляемого аппарата, вторая ---
положению цели в $t_f$. После вычисления скоростей дуги
в инерциальных осях определяются импульсы
\begin{equation}
 \Delta\vect V_{1,I}=\vect V_1-\vect V_C(t_i^-),\qquad
 \Delta\vect V_{2,I}=\vect V_{\rm tar}(t_f)-\vect V_2.
 \label{eq:lambert-rendezvous-control}
\end{equation}
Для получения компонент в ОСК импульсы проецируют
в предманёвренные базисы. Это решение для выбранных
времени и ветви. В центральном поле оно удовлетворяет
нелинейному конечному условию; при возмущениях
используется как начальное приближение.

\FloatBarrier
\section{Относительное движение и модель CW}
\label{sec:cw}

\subsection{Относительная система координат}

Относительные координаты позволяют рассматривать движение
между аппаратами в масштабе их расстояния.
Обозначим положение цели через $\vect R_0(t)$,
положение управляемого аппарата --- через $\vect R_C(t)$.
Базис RTN цели определён формулами \eqref{eq:RTN}; матрица
$Q_0$ содержит его векторы в столбцах. Относительное положение
в локальных компонентах равно
\begin{equation}
 \vect\rho=Q_0^\top(\vect R_C-\vect R_0)
           =(x,y,z)^\top.
 \label{eq:relative-position}
\end{equation}
Для кеплеровой цели базис поворачивается в неподвижной
плоскости с угловой скоростью
$\vect\Omega_r=(0,0,\omega_r)^\top$,
где $\omega_r=\dot\nu=h/R_0^2$.
Индекс отличает эту скорость от аргумента перицентра $\omega$.

Производная координат во вращающемся базисе связана
с разностью инерциальных скоростей соотношением
\begin{equation}
 Q_0^\top(\vect V_C-\vect V_0)
 =\dot{\vect\rho}+\vect\Omega_r\times\vect\rho.
 \label{eq:relative-velocity}
\end{equation}
Следовательно,
\begin{equation}
 \delta V_r=\dot x-\omega_r y,\qquad
 \delta V_t=\dot y+\omega_r x,\qquad
 \delta V_n=\dot z.
 \label{eq:velocity-components-relative}
\end{equation}
При $\dot{\vect\rho}=0$ положение в RTN постоянно,
но разность инерциальных скоростей равна
$\vect\Omega_r\times\vect\rho$.

\begin{explanation}{Согласованное начальное состояние}
Если задано $\vect\rho_i,\dot{\vect\rho}_i$, инерциальное
состояние второго аппарата строится по формулам
\[
 \vect R_{Ci}=\vect R_{0i}+Q_{0i}\vect\rho_i,\qquad
 \vect V_{Ci}=\vect V_{0i}
 +Q_{0i}(\dot{\vect\rho}_i+
                \vect\Omega_{ri}\times\vect\rho_i).
\]
\end{explanation}

\subsection{Точные нелинейные уравнения относительно кеплеровой цели}

Для обоих аппаратов сначала примем одну и ту же центральную
гравитационную модель. Вычтем их инерциальные уравнения.
Дважды используя транспортную формулу вращающегося базиса,
получим
\begin{equation}
 \begin{aligned}
 \ddot{\vect\rho}={}&-\mu\left[
 \frac{R_0\vect e_x+\vect\rho}
      {\norm{R_0\vect e_x+\vect\rho}^3}
       -\frac{\vect e_x}{R_0^2}\right]
 -2\vect\Omega_r\times\dot{\vect\rho}
 -\dot{\vect\Omega}_r\times\vect\rho\\
 &-\vect\Omega_r\times(\vect\Omega_r\times\vect\rho)
 +\vect u.
 \end{aligned}
 \label{eq:exact-relative}
\end{equation}
Здесь $\vect e_x=(1,0,0)^\top$ --- компоненты радиального
направления в локальном базисе, $\vect u$ --- разность
добавочных ускорений в том же базисе. Для пассивной цели
и управляемого аппарата это управляющее ускорение аппарата.
Уравнение сохраняет нелинейность по относительному расстоянию.

Положим $d=\sqrt{(R_0+x)^2+y^2+z^2}$. В компонентах
\begin{equation}
 \begin{aligned}
 \ddot x-2\omega_r\dot y-\dot\omega_r y-\omega_r^2x
   &=-\frac{\mu(R_0+x)}{d^3}+\frac{\mu}{R_0^2}+u_x,\\
 \ddot y+2\omega_r\dot x+\dot\omega_r x-\omega_r^2y
   &=-\frac{\mu y}{d^3}+u_y,\\
 \ddot z&=-\frac{\mu z}{d^3}+u_z.
 \end{aligned}
 \label{eq:exact-relative-components}
\end{equation}
Здесь учтены кориолисовы, эйлеровы и центробежные члены.

\subsection{Линеаризация вокруг круговой орбиты}

Пусть цель движется по кругу радиуса $R_0=a_0$.
Тогда $\omega_r=n_0=\sqrt{\mu/a_0^3}$ постоянно,
$\dot\omega_r=0$. При $\norm{\vect\rho}/a_0\ll1$
разность гравитационных ускорений в первом порядке равна
\begin{equation}
 -\mu\left[
 \frac{a_0\vect e_x+\vect\rho}{\norm{a_0\vect e_x+\vect\rho}^3}
 -\frac{\vect e_x}{a_0^2}\right]
 =n_0^2(2x,-y,-z)^\top
    +O\left(\frac{\mu\norm{\vect\rho}^2}{a_0^4}\right).
 \label{eq:gravity-gradient-circular}
\end{equation}
Подставив эту разность в \eqref{eq:exact-relative-components},
получим уравнения Клохесси--Уилтшира (CW):
\begin{equation}
 \boxed{\begin{aligned}
  \ddot x-2n_0\dot y-3n_0^2x&=u_x,\\
  \ddot y+2n_0\dot x&=u_y,\\
  \ddot z+n_0^2z&=u_z.
 \end{aligned}}
 \label{eq:CW}
\end{equation}
Уравнения CW \cite{cw} применимы для круговой кеплеровой
опорной орбиты и малых относительных расстояний.
Декартовы координаты определены при $e=0$,
хотя классические элементы в этом случае вырождены.

\subsection{Свободное движение и дрейф}

При $\vect u=0$ интегрирование второго уравнения даёт
\begin{equation}
 \dot y+2n_0x=C,\qquad C=\dot y_i+2n_0x_i.
 \label{eq:CW-invariant}
\end{equation}
Подставив $\dot y=C-2n_0x$ в первое уравнение, получаем
\[
 \ddot x+n_0^2x=2n_0C.
\]
Его общее решение и последующее интегрирование $\dot y$ дают
\begin{equation}
 \begin{aligned}
 x(t)&=A\cos(n_0t)+B\sin(n_0t)+\frac{2C}{n_0},\\
 y(t)&=-2A\sin(n_0t)+2B\cos(n_0t)-3Ct+D,\\
 z(t)&=E\cos(n_0t)+F\sin(n_0t).
 \end{aligned}
 \label{eq:CW-modes}
\end{equation}
Начало времени совмещено с $t_i$.
Постоянная $C$ одновременно задаёт смещение центра радиальных
колебаний $x_c=2C/n_0$ и продольный дрейф $-3Ct$.
Отсюда
\begin{equation}
 \dot y_{\rm secular}=-\frac32n_0x_c.
 \label{eq:CW-drift}
\end{equation}
При положительном радиальном смещении центра движение
в среднем отстаёт от цели. Дрейф соответствует зависимости среднего движения
от большой полуоси в \eqref{eq:phase-variation}.

Для отсутствия секулярного дрейфа необходимо
\begin{equation}
 C=0\quad\Longleftrightarrow\quad
 \dot y_i=-2n_0x_i.
 \label{eq:CW-no-drift}
\end{equation}
При этом плоская траектория является эллипсом с отношением
полуосей $2:1$, а $D$ смещает его центр вдоль орбиты.
\coursefigure{cw-modes}{Два свободных движения CW при радиусе
опорной круговой орбиты 7000~км. Слева $x_i=100$~м,
$\dot x_i=0$, $\dot y_i=-2n_0x_i$; справа изменена
начальная продольная производная, поэтому появляется дрейф.
Горизонтальная ось рисунка --- $y$, вертикальная --- $x$.}

\subsection{Решение через начальное состояние}

Положим $c=\cos(n_0t)$, $s=\sin(n_0t)$.
После определения постоянных в \eqref{eq:CW-modes} получаем
\begin{equation}
 \begin{aligned}
 x&=(4-3c)x_i+\frac{s}{n_0}\dot x_i
                        +\frac{2(1-c)}{n_0}\dot y_i,\\
 y&=6(s-n_0t)x_i+y_i-\frac{2(1-c)}{n_0}\dot x_i
                        +\frac{4s-3n_0t}{n_0}\dot y_i,\\
 z&=cz_i+\frac{s}{n_0}\dot z_i,\\
 \dot x&=3n_0s x_i+c\dot x_i+2s\dot y_i,\\
 \dot y&=6n_0(c-1)x_i-2s\dot x_i+(4c-3)\dot y_i,\\
 \dot z&=-n_0s z_i+c\dot z_i.
 \end{aligned}
 \label{eq:CW-initial-solution}
\end{equation}
При отсутствии дрейфа все три координаты являются
линейными комбинациями $\sin n_0t$, $\cos n_0t$ и постоянного
сдвига. Пространственная траектория представляет эллипс.

\coursefigure{cw-spatial}{Пространственное свободное движение CW
без дрейфа: $x=100\cos(n_0t)$~м,
$y=-200\sin(n_0t)$~м,
$z=120\sin(n_0t+0.6)$~м. Справа показаны те же
компоненты во времени; общая частота обеспечивает замкнутость.}

\subsection{Управление и линейное конечное условие в RTN}

Пусть
$\vect X=(x,y,z,\dot x,\dot y,\dot z)^\top$.
Уравнения \eqref{eq:CW} записываются как
\[
 \dot{\vect X}=A_{\rm CW}\vect X+G\vect u,
 \qquad G=\begin{pmatrix}0_{3\times3}\\I_3\end{pmatrix}.
\]
Мгновенный импульс, выраженный в RTN цели, даёт
$\vect X^+=\vect X^-+G\Delta\vect V$.
Решение с $N$ импульсами имеет вид
\begin{equation}
 \vect X_f=\Phi_{\rm CW}(t_f,t_i)\vect X_i
       +\sum_{j=1}^N\Phi_{\rm CW}(t_f,t_j)G\Delta\vect V_j.
 \label{eq:CW-guidance}
\end{equation}
Условие встречи $\vect X_f=0$ обеспечивает совпадение
положения и инерциальной скорости.
Структура конечного условия аналогична \eqref{eq:linear-boundary},
но матрицы состояния относятся к декартовым координатам RTN.

\FloatBarrier
\section{Эллиптическая орбита: TH и решение Яманаки--Анкерсена}
\label{sec:th}

\subsection{Уравнения движения}

На эллиптической опорной орбите расстояние до центрального тела
и скорость вращения RTN меняются в течение витка:
\begin{equation}
 R_0=\frac{p_0}{1+e_0\cos\nu_0},\qquad
 \omega_r=\frac{h_0}{R_0^2},\qquad
 \dot\omega_r=-2\frac{\dot R_0}{R_0}\omega_r.
 \label{eq:ellipse-reference-rates}
\end{equation}
В TH коэффициенты не заменяются постоянным средним движением.
Линеаризуем гравитационную разность
\eqref{eq:exact-relative} по $\vect\rho$, не разлагая
коэффициенты по $e_0$. Получим временную форму
уравнений Чаунера--Хемпеля (TH):
\begin{equation}
 \boxed{\begin{aligned}
 \ddot x&=\left(\omega_r^2+2\frac{\mu}{R_0^3}\right)x
               +2\omega_r\dot y+\dot\omega_r y+u_x,\\
 \ddot y&=\left(\omega_r^2-\frac{\mu}{R_0^3}\right)y
               -2\omega_r\dot x-\dot\omega_r x+u_y,\\
 \ddot z&=-\frac{\mu}{R_0^3}z+u_z.
 \end{aligned}}
 \label{eq:TH-time}
\end{equation}
При $e_0\to0$ имеем
$\omega_r^2=\mu/R_0^3=n_0^2$, $\dot\omega_r=0$,
и система переходит в CW с теми же знаками и осями.
Модель сохраняет первый порядок по относительному расстоянию
при любом $0\le e_0<1$. Пассивные возмущения здесь не учитываются.

\subsection{Переход к истинной аномалии и масштабирование}

Для пассивной системы выберем независимой переменной
$\nu_0$, а штрихом обозначим производную по ней.
Положим
\begin{equation}
 \sigma=1+e_0\cos\nu_0,\qquad
 \widetilde{\vect\rho}=\sigma\vect\rho.
 \label{eq:TH-scaling}
\end{equation}
Используются соотношения
\[
 \frac{\dd}{\dd t}=\omega_r\frac{\dd}{\dd\nu_0},\qquad
 \ddot a=\omega_r^2a''+\dot\omega_r a'
\]
для любой скалярной функции $a$.
Дифференцирование произведения в \eqref{eq:TH-scaling}
позволяет убрать первые производные, возникающие
из переменной угловой скорости.
После подстановки в \eqref{eq:TH-time} получаем
\begin{equation}
 \begin{aligned}
  \widetilde x''-2\widetilde y'
                    -\frac{3}{\sigma}\widetilde x&=0,\\
  \widetilde y''+2\widetilde x'&=0,\\
  \widetilde z''+\widetilde z&=0.
 \end{aligned}
 \label{eq:TH-anomaly}
\end{equation}
Производные в \eqref{eq:TH-anomaly} берутся по аномалии,
а координаты масштабированы множителем $\sigma$.

Из второй строки следует
$\widetilde y'+2\widetilde x=K$, где $K$ постоянно.
Первая строка тогда принимает вид
\begin{equation}
 \widetilde x''+
     \left(4-\frac3\sigma\right)\widetilde x=2K.
 \label{eq:TH-scalar-radial}
\end{equation}
Одна однородная функция решения равна
$\sigma\sin\nu_0$. Остальные функции находятся понижением порядка и квадратурами.

\subsection{Матрица перехода Яманаки--Анкерсена}

Яманака и Анкерсен \cite{ya} получили удобную аналитическую
матрицу перехода той же линейной эллиптической системы.
YA является аналитическим решением TH, регулярным при $e_0=0$.
Важная для него интегральная величина
\begin{equation}
 \mathcal J(\nu_0,\nu_i)
 =\int_{\nu_i}^{\nu_0}\frac{\dd\eta}{(1+e_0\cos\eta)^2}
 =\sqrt{\frac{\mu}{p_0^3}}(t-t_i)
 \label{eq:YA-J}
\end{equation}
сохраняет информацию о прошедшем времени и числе оборотов.
Обозначим фундаментальную матрицу системы
\eqref{eq:TH-anomaly} через $Y(\nu_0)$.
Она состоит из шести независимых аналитических решений.
Для связи с физическим состоянием
$\vect X=(\vect\rho,\dot{\vect\rho})^\top$ используется
преобразование
\begin{equation}
 \widetilde{\vect X}=
 \begin{pmatrix}\widetilde{\vect\rho}\\
                 \widetilde{\vect\rho}'\end{pmatrix}
 =\mathcal T(\nu_0)\vect X,
 \qquad
 \mathcal T=
 \begin{pmatrix}
 \sigma I_3&0\\
 \sigma'I_3&(\sigma/\omega_r)I_3
 \end{pmatrix}.
 \label{eq:YA-transform}
\end{equation}
Тогда матрица перехода в координатах и производных
по физическому времени имеет вид
\begin{equation}
 \Phi_{\rm TH}(t,t_i)=
 \mathcal T(\nu_0)^{-1}Y(\nu_0)Y(\nu_i)^{-1}
                              \mathcal T(\nu_i).
 \label{eq:YA-STM-concept}
\end{equation}
Преобразование \eqref{eq:YA-transform} учитывает и масштаб
координат, и замену времени аномалией.

В первоисточнике YA оси выбраны иначе: $X$ трансверсален,
$Y$ противоположен нормали, $Z$ направлен в надир.
Наши координаты связаны с ними как
$x=-Z$, $y=X$, $z=-Y$. После этой замены временная система совпадает с
\eqref{eq:TH-time}.

\subsection{Сравнение CW и TH}

Возьмём цель на эллипсе $a_0=10000$~км, $e_0=0.2$,
старт в перицентре. В её RTN зададим
\[
 \vect\rho_i=(0,-100,10)^\top\ \text{м},\qquad
 \dot{\vect\rho}_i=(0,0,0)^\top\ \text{м/с}.
\]
Точное состояние второго аппарата строится по
\eqref{eq:relative-velocity}. Затем независимо прогнозируются
два нелинейных двухтельных движения. Их разность проецируется
на текущую RTN реальной эллиптической цели.
Коэффициенты TH берутся по той же цели;
в приближении CW явно используется $n_0=\sqrt{\mu/a_0^3}$.

На интервале одного периода, около $9952$~с,
максимальная ошибка CW составляет примерно $50.27$~м,
а TH --- около $0.050$~м. В конце периода ошибки обеих
моделей близки к $0.050$~м. Максимальная и конечная ошибки характеризуют
разные участки движения.

\coursefigure{relative-elliptic}{Сравнение относительно одной
эллиптической цели: точное двухтельное движение, TH и CW.
В этом специальном примере линейная эллиптическая проекция
имеет $x=0$, а $y$ меняется пропорционально радиусу цели;
CW даёт неподвижную точку в проекции. Справа видна
ошибка на всём интервале, а не только в конце витка.}

По TH в этом примере $\widetilde y$ постоянна,
поэтому $y$ пропорциональна $R_0$.
Радиальная составляющая точного решения имеет второй порядок.

\FloatBarrier
\section{Околокруговая орбита с учётом \texorpdfstring{$J_2$}{J2}: модель Швайгарта--Седвика}
\label{sec:ss}

\subsection{Влияние второй зональной гармоники}

Вторая зональная гармоника вызывает прецессию плоскости
орбиты и изменение фазовых частот. Для длительного относительного
движения эти эффекты могут быть существенны.

Обозначим экваториальный радиус Земли через $R_E$,
вторую зональную гармонику --- через $J_2$.
В инерциальных экваториальных осях добавочное ускорение равно
\begin{equation}
 \vect a_{J_2}=\frac{3J_2\mu R_E^2}{2R^5}
 \begin{pmatrix}
 X(5Z^2/R^2-1)\\
 Y(5Z^2/R^2-1)\\
 Z(5Z^2/R^2-3)
 \end{pmatrix}.
 \label{eq:J2-acceleration}
\end{equation}
Здесь $X,Y,Z$ --- инерциальные координаты аппарата,
$R=\sqrt{X^2+Y^2+Z^2}$. Модель SS \cite{ss} описывает относительное движение близких
околокруговых аппаратов с основными эффектами $J_2$.
Полный прогноз каждого аппарата использует
\eqref{eq:J2-acceleration} без усреднения.

В первом порядке усреднённые скорости классических элементов
дают знакомую картину прецессии:
\begin{equation}
 \begin{aligned}
 \dot{\overline\Omega}
 &=-\frac32J_2n\left(\frac{R_E}{p}\right)^2\cos i,\\
 \dot{\overline\omega}
 &=\frac34J_2n\left(\frac{R_E}{p}\right)^2(5\cos^2i-1),\\
 \dot{\overline M}
 &=n+\frac34J_2n\left(\frac{R_E}{p}\right)^2
                   \sqrt{1-e^2}(3\cos^2i-1).
 \end{aligned}
 \label{eq:J2-secular}
\end{equation}
Черта обозначает усреднённое описание.
Оскулирующие элементы дополнительно имеют короткопериодические
изменения. Поэтому усреднённые элементы нельзя без преобразования
отождествлять с мгновенными элементами, найденными из одного
инерциального состояния.

\subsection{Опорная орбита и усреднение}

В SS вводится опорная орбита постоянного радиуса $r_{\rm ref}$,
скорректированная по среднему движению и ориентации
с учётом $J_2$. Градиент возмущающего ускорения усредняется.
Для выбранного наклонения $i_{\rm ref}$
вводятся безразмерные коэффициенты
\begin{equation}
 s=\frac{3J_2R_E^2}{8r_{\rm ref}^2}
                   (1+3\cos2i_{\rm ref}),\qquad
 c=\sqrt{1+s},\qquad n=\sqrt{\frac{\mu}{r_{\rm ref}^3}}.
 \label{eq:SS-constants}
\end{equation}
В локальных радиальном, трансверсальном и нормальном
направлениях усреднённый градиент $J_2$ имеет вид
\begin{equation}
 \left\langle\frac{\partial\vect a_{J_2}}
                         {\partial\vect R}\right\rangle
 =n^2\diag(4s,-s,-3s).
 \label{eq:SS-averaged-gradient}
\end{equation}
Частота вращения орбитального базиса равна $nc$.
Короткопериодические изменения градиента при усреднении
отбрасываются. Для наклонённой орбиты SS также корректирует
ориентацию опорной плоскости и нормальное движение.
При сравнении с точным прогнозом состояния переводятся
в одну систему координат.

\subsection{Плоская часть и её частота}

Для разности координат двух близких аппаратов плоская
часть уравнения (45) работы \cite{ss} имеет вид
\begin{equation}
 \begin{aligned}
 \ddot x-2nc\dot y-(5c^2-2)n^2x&=0,\\
 \ddot y+2nc\dot x&=0.
 \end{aligned}
 \label{eq:SS-planar}
\end{equation}
Здесь $x$ радиален, $y$ трансверсален в скорректированном базисе.
Коэффициент первой строки можно проверить непосредственно.
Радиальный гравитационный градиент центрального поля равен
$2n^2$, усреднённая добавка --- $4sn^2$, центробежный
член --- $c^2n^2$. Их сумма
$2+4s+c^2=5c^2-2$ даёт именно указанный коэффициент.
В трансверсальном направлении
$-n^2-sn^2+c^2n^2=0$.

Интегрируя вторую строку, положим
\begin{equation}
 \dot y+2ncx=C_s.
 \label{eq:SS-invariant}
\end{equation}
Тогда
\begin{equation}
 \ddot x+\beta^2x=2ncC_s,\qquad
 \beta=n\sqrt{1-s}.
 \label{eq:SS-planar-frequency}
\end{equation}
Таким образом, частота плоских колебаний $\beta$
отличается от скорости вращения базиса $nc$.
При $J_2\to0$ имеем $s\to0$, $c\to1$, $\beta\to n$,
и плоская система переходит в CW.

Для $C_s=0$ решение можно записать
\begin{equation}
 \begin{aligned}
 x&=A\cos\beta t+B\sin\beta t,\\
 y&=-\frac{2nc}{\beta}A\sin\beta t
          +\frac{2nc}{\beta}B\cos\beta t+D.
 \end{aligned}
 \label{eq:SS-bounded-planar}
\end{equation}
Отношение полуосей равно $2c/\sqrt{1-s}$, а не строго двум.
Чтобы убрать дрейф, требуется $\dot y_i=-2ncx_i$.
Чтобы дополнительно убрать постоянный продольный сдвиг,
нужно $\dot x_i=\beta^2y_i/(2nc)$.
Оба условия следуют из общего решения.

При $C_s\ne0$ центр радиальных колебаний равен
$x_c=2ncC_s/\beta^2$, а скорость секулярного изменения $y$
\begin{equation}
 \dot y_{\rm secular}=-\frac{n(5c^2-2)}{2c}\,x_c.
 \label{eq:SS-drift}
\end{equation}
При $s=0$ формулы переходят в CW.

\subsection{Нормальная компонента и геометрия прецессии}

Если ограничиться одним усреднённым градиентом
\eqref{eq:SS-averaged-gradient}, получится простое уравнение
$\ddot z+n^2(1+3s)z=0$. Это промежуточное приближение:
в нём не полностью восстановлена геометрия движения
пересечения двух прецессирующих плоскостей.
Полная нормальная часть уравнения (45) SS имеет вид
\begin{equation}
 \ddot z+q^2z=2lq\cos(qt+\varphi).
 \label{eq:SS-cross}
\end{equation}
Параметры $q,l,\varphi$ определяются взаимной геометрией
орбитальных плоскостей и скоростями их прецессии.
Пусть $\Psi(t)$ --- малый угол между нормалями двух
плоскостей. Его величина определяется равенством
\begin{equation}
 \cos\Psi=\cos i_1\cos i_2+
                    \sin i_1\sin i_2\cos(\Omega_1-\Omega_2).
 \label{eq:mutual-plane-angle}
\end{equation}
При $\dot\Omega_1\ne\dot\Omega_2$ меняется линия пересечения плоскостей.
Относительное нормальное положение можно представить
как
\begin{equation}
 z=A(t)\sin\chi(t),\qquad
 A(t)\simeq r_{\rm ref}\sin\Psi(t).
 \label{eq:SS-amplitude-geometry}
\end{equation}
Угол $\chi$ отсчитывается от ориентированной линии
пересечения в плоскости соответствующей орбиты.
Положительный знак нормали и ориентация линии должны
быть выбраны согласованно.

Локальная аппроксимация
$A(t)=m+lt$, $\chi(t)=qt+\varphi$ даёт
\begin{equation}
 z=(m+lt)\sin(qt+\varphi),\qquad
 l=\dot A(0),\quad q=\dot\chi(0).
 \label{eq:SS-cross-solution}
\end{equation}
Двойное дифференцирование немедленно приводит
к \eqref{eq:SS-cross}. Начальные условия удовлетворяют
\[
 m\sin\varphi=z_i,\qquad
 l\sin\varphi+qm\cos\varphi=\dot z_i.
\]
Для вычисления угла $\chi$ можно построить единичные
нормали $\vect n_1,\vect n_2$, ориентированное пересечение
$\vect k=(\vect n_2\times\vect n_1)/\sin\Psi$
и вектор $\vect l_1=\vect n_1\times\vect k$.
Тогда $\chi=\atanTwo(\widehat{\vect R}_1\cdot\vect l_1,
\widehat{\vect R}_1\cdot\vect k)$, а его производная
включает движение аппарата и движение пересечения.
При совпадающих плоскостях такое построение имеет
координатную особенность; необходимо переходить к пределу
или работать непосредственно с векторами, не деля
на почти нулевой $\sin\Psi$.

Различие частот $\beta,q$ и изменение амплитуды
поворачивают пространственную относительную фигуру.

\coursefigure{j2-frequencies}{Геометрический пример
дифференциальной прецессии: изменение нормального размаха
вычислено по усреднённым классическим скоростям узлов
и углу между плоскостями. Рисунок иллюстрирует механизм;
он не является проверкой всех коэффициентов полной модели SS.}

\subsection{Пример на экваториальной орбите}

На экваторе $i_{\rm ref}=0$ и
$s=\tfrac32J_2(R_E/r_{\rm ref})^2$.
В полном поле центрального притяжения и $J_2$
существует точная экваториальная круговая орбита
с частотой $\omega_{\rm ref}=n\sqrt{1+s}$.
Плоское сравнение выполняется в RTN этой орбиты.

На рисунке~\ref{fig:j2-relative} взято
$r_{\rm ref}=7000$~км, $x_i=100$~м, $y_i=0$,
$\dot x_i=0$, $\dot y_i=-2\omega_{\rm ref}x_i$.
Все модели используют одни и те же начальные положения
и скорости обоих аппаратов. В модели CW положено
$n=\omega_{\rm ref}$, поэтому обе линейные модели используют
одинаковую частоту вращения системы координат.
Остающееся различие связано с коэффициентами относительной
динамики и частотой радиальных колебаний.

\coursefigure{j2-relative}{Плоский экваториальный пример
с $J_2$ на 30 витках. Сравниваются два полных нелинейных
движения, CW с согласованной частотой вращения цели
и плоская часть SS. Это проверка одного частного случая;
нормальная динамика и произвольные наклонения здесь
не проверяются.}

В этом расчёте максимальная ошибка положения CW составляет
около $50.8$~м, а плоской SS --- около $0.41$~м.
Уточнение допусков полного прогноза меняет сравнение
примерно на $10^{-5}$~м. Здесь проверена только
экваториальная плоская задача.

\subsection{Области применения}

CW используется около круговой кеплеровой орбиты,
TH/YA --- около эллиптической кеплеровой орбиты,
SS --- при существенных эффектах $J_2$ в окрестности круговой орбиты.
Для более общей модели сил матрицу перехода получают
численным интегрированием вариационных уравнений.
Относительные орбитальные элементы являются другим
представлением состояния и сами по себе не добавляют возмущений.
Во всех линейных моделях требуется малое удаление от опорной траектории.

\FloatBarrier
\section{Сквозной расчёт встречи на близких орбитах}
\label{sec:example}

\subsection{Исходные данные}

Рассмотрим двухтельную задачу вокруг Земли
с $\mu=398600.4418$~км$^3$/с$^2$.
Начальные элементы цели и отклонения элементов
управляемого аппарата заданы следующими значениями:
\begin{equation}
 \begin{aligned}
 \ExampleInitialData
 \end{aligned}
 \label{eq:example-data}
\end{equation}
Индекс 0 относится к цели, $\delta$ означает «аппарат минус цель».
Начальное расстояние равно $\ExampleInitialDistance$.
Два импульса прикладываются при $t_1=0$ и $t_2=t_f=1200$~с.
Конечное состояние задаётся после всех манёвров.
Классические элементы регулярны: $e>0$, $\sin i\ne0$.

\subsection{Решение в линейной элементной модели}

Независимо рассчитываем движение цели и движение аппарата
без управления. По этим траекториям вычисляются
$\Delta\vect y_f=\vect y_{\rm tar}(t_f)-\vect y_0(t_f)$,
$B_1,B_2$ и
\[
 \mathcal M=
 [\Phi(t_f,0)B_1\quad \Phi(t_f,t_f)B_2]
 =[\Phi(t_f,0)B_1\quad B_2].
\]
Квадратная система невырождена. Её решение приведено
в таблице в м/с.

\begin{table}[htbp]
\centering
\caption{Линейное управление в предманёвренных ОСК}
\begin{tabular}{@{}lrrr@{}}
\toprule
Импульс&$\Delta V_r$&$\Delta V_t$&$\Delta V_n$\\\midrule
\ExampleLinearComponents
\bottomrule
\end{tabular}
\end{table}

Характеристическая скорость равна $\ExampleLinearCost$.

\subsection{Решение через задачу Ламберта}

Начальное положение аппарата и положение цели
через 1200~с подставляются в задачу
\eqref{eq:lambert-boundary}. Рассматривается
одна дуга без полных оборотов в направлении движения,
близком к начальному. Угол перелёта равен
$\ExampleTransferAngle$.
Скалярное уравнение времени решается по $p$;
полученный фокальный параметр равен $\ExampleLambertP$.
По \eqref{eq:lambert-velocities} восстанавливаются
скорости, затем по
\eqref{eq:lambert-rendezvous-control} --- импульсы.

\begin{table}[htbp]
\centering
\caption{Управление по задаче Ламберта в тех же единицах}
\begin{tabular}{@{}lrrr@{}}
\toprule
Импульс&$\Delta V_r$&$\Delta V_t$&$\Delta V_n$\\\midrule
\ExampleLambertComponents
\bottomrule
\end{tabular}
\end{table}

Характеристическая скорость равна $\ExampleLambertCost$.

\coursefigure{rendezvous-orbits}{Сквозной пример в инерциальном
пространстве: траектория цели и два управляемых перелёта,
рассчитанных линейной моделью и методом Ламберта.
Справа показано начальное разнесение аппаратов.
На орбитальном масштабе близкие траектории почти накладываются.}

\subsection{Относительная траектория}

Относительное положение вычисляется в RTN цели:
$\vect\rho(t)=Q_0(t)^\top(\vect R_C(t)-\vect R_0(t))$.

\coursefigure{rendezvous-relative}{Тот же расчёт в системе цели.
Начальная точка и конечное время общие для сравниваемых
управлений. Относительные проекции и временные компоненты
показывают геометрию движения, скрытую на инерциальном
масштабе.}

\subsection{Нелинейная проверка управления}

Каждое управление применяется к одному исходному
состоянию. Положение и скорость преобразуются
в инерциальные координаты, выполняется точный
скачок \eqref{eq:exact-impulse}, между манёврами
прогнозируется двухтельное движение. После
заключительного манёвра определяются
\begin{equation}
 \varepsilon_R=\norm{\vect R_C(t_f)-\vect R_0(t_f)},\qquad
 \varepsilon_V=\norm{\vect V_C(t_f)-\vect V_0(t_f)}.
 \label{eq:physical-terminal-errors}
\end{equation}
Эти две величины имеют непосредственный физический
смысл и не зависят от выбора орбитальных элементов.

\begin{table}[htbp]
\centering
\caption{Проверка управлений исходной двухтельной динамикой}
\begin{tabular}{@{}lrrr@{}}
\toprule
Решение&$\sum\|\Delta\vect V_j\|$, м/с&
$\varepsilon_R$, м&$\varepsilon_V$, м/с\\\midrule
\ExampleComparisonRows
\bottomrule
\end{tabular}
\end{table}

Линейное управление даёт $\varepsilon_R=0.492$~м
и $\varepsilon_V\simeq7.4\cdot10^{-4}$~м/с.
Для Ламберта и уточнённого решения невязки ниже
$10^{-6}$~м и $10^{-9}$~м/с показаны верхними границами.
Независимое интегрирование инерциальных уравнений
согласуется с аналитическим прогнозом примерно до
$2\cdot10^{-7}$~м и $5\cdot10^{-10}$~м/с.
Это вычислительная проверка двухтельной модели;
ошибки навигации и исполнения импульсов здесь не учитываются.

\FloatBarrier
\section{Нелинейное уточнение линейного решения}
\label{sec:correction}

\subsection{Конечная невязка и матрица чувствительности}

Число импульсов, моменты их приложения и $t_f$ фиксируются.
Уточняются компоненты импульсов в предманёвренных ОСК.
Обозначим через $\vect y_f(\Delta\vect V)$ конечное состояние,
полученное с точными скачками скорости и нелинейным движением
между манёврами. Состояние записывается в классических элементах
$(a,e,\omega,i,\Omega,M)$.

По аналогии с \cite[\S~4.2]{chiniaev} на итерации $k$ вводим
\begin{equation}
 \delta\vect y_f^{(k)}=\vect y_t-\vect y_f(\Delta\vect V^{(k)}),
 \qquad
 \mathcal M^{(k)}=
 \left.\frac{\partial\vect y_f}{\partial\Delta\vect V}
                         \right|_{\Delta\vect V^{(k)}}.
 \label{eq:active-Jacobian}
\end{equation}
Линеаризация конечного состояния даёт уравнение поправки
\begin{equation}
 \mathcal M^{(k)}\delta\vect V^{(k)}
                  =\delta\vect y_f^{(k)}.
 \label{eq:Newton-step}
\end{equation}
В статье фиксируются истинные долготы манёвров;
здесь фиксируются моменты времени. При задании манёвров
по долготам требуется дополнительно учитывать изменение
времени достижения каждой долготы.

В точке $\Delta\vect V=0$ матрица чувствительности сводится
к блокам $\Phi B$ на траектории без управления.
При ненулевых импульсах производная вычисляется заново:
изменяются орбиты между манёврами и предманёвренные базисы.
Использование одной неизменной матрицы $\mathcal M$
соответствует модифицированной итерации с замороженными
производными, а не точному методу Ньютона.

\subsection{Вычисление поправки}

Для небольшого числа манёвров столбцы матрицы удобно
вычислять центральными разностями:
\begin{equation}
 (\mathcal M^{(k)})_{:,\ell}\approx
 \frac{\vect y_f(\Delta\vect V^{(k)}+h_\ell\vect e_\ell)
       -\vect y_f(\Delta\vect V^{(k)}-h_\ell\vect e_\ell)}{2h_\ell}.
 \label{eq:finite-difference-J}
\end{equation}
Здесь $h_\ell$ --- приращение компоненты импульса,
$\vect e_\ell$ --- единичный столбец пространства управления.
В каждом прогнозе заново строятся ОСК всех манёвров.
Шаг проверяют повторным вычислением при $h_\ell/2$.
Слишком малый шаг усиливает округление, слишком большой ---
ошибку разностной производной. В примере взято
$h_\ell=10^{-3}$~м/с.

Углы $\omega,\Omega,M$ в конечной невязке и разностях
продолжаются около одной локальной ветви.
Разрыв записи на $2\pi$ не должен попадать в производную.
Малые $e$ и $\sin i$ требуют отдельной проверки
регулярности классических координат.

Для двух импульсов система \eqref{eq:Newton-step} квадратная.
При невырожденной матрице поправка единственна.
Для $N>2$ при полном строчном ранге минимальная поправка
по евклидовой норме определяется задачей
\begin{equation}
 \min_{\delta\vect V^{(k)}}\sum_{j=1}^N
                \norm{\delta\vect V_j^{(k)}}^2,
 \qquad
 \mathcal M^{(k)}\delta\vect V^{(k)}=\delta\vect y_f^{(k)}.
 \label{eq:correction-minnorm-problem}
\end{equation}
Её решение
\begin{equation}
 \delta\vect V^{(k)}=(\mathcal M^{(k)})^\top
 \left[\mathcal M^{(k)}(\mathcal M^{(k)})^\top\right]^{-1}
                                      \delta\vect y_f^{(k)}
 \label{eq:correction-minnorm}
\end{equation}
минимизирует норму поправки к управлению. Сумма норм
полных импульсов в этой подзадаче не минимизируется.
При активных ограничениях
$\norm{\Delta\vect V_j^{(k)}+\delta\vect V_j^{(k)}}
\le\Delta V_{\max,j}$ их включают в подзадачу.
При недостаточном ранге полная компенсация невязки
может быть невозможна.

\subsection{Обновление управления и остановка}

После вычисления поправки выполняется
\begin{equation}
 \Delta\vect V^{(k+1)}=\Delta\vect V^{(k)}
                   +\alpha_k\delta\vect V^{(k)},
 \qquad 0<\alpha_k\le1.
 \label{eq:damped-update}
\end{equation}
Если полный шаг увеличивает масштабированную конечную
невязку, $\alpha_k$ уменьшают.
Для численного решения строки классического состояния
масштабируют; для остановки проверяют отдельно
ошибки положения и скорости:
\begin{equation}
 \varepsilon_R\le\varepsilon_{R,\rm req},\qquad
 \varepsilon_V\le\varepsilon_{V,\rm req}.
 \label{eq:physical-stopping}
\end{equation}
Точность прогноза должна позволять различать эти допуски.
Малый шаг при большой невязке не является выполнением
конечного условия.

\subsection{Уточнение расчётного примера}

Начальные импульсы получены линейной элементной моделью;
их моменты 0 и 1200~с фиксированы.
На каждой итерации заново вычисляются шесть столбцов
\eqref{eq:finite-difference-J} и решается
\eqref{eq:Newton-step}. В этом случае принимается полный шаг.

\begin{table}[htbp]\centering
\caption{Уточнение по классической конечной невязке}
\begin{tabular}{@{}rrrr@{}}
\toprule
Итерация&$\varepsilon_R$, м&$\varepsilon_V$, м/с&
$\sum\|\Delta\vect V_j\|$, м/с\\\midrule
\ExampleCorrectionRows
\bottomrule
\end{tabular}
\end{table}

Одна поправка уменьшает ошибку положения с $0.492$~м
до $1.1\cdot10^{-6}$~м. Характеристическая скорость меняется
с $4.40864$ до $4.40860$~м/с.

\coursefigure{rendezvous-convergence}{Конечные ошибки положения
и скорости при обновлении матрицы чувствительности.
Поправка вычисляется по классическим элементам.
Значения ниже уровня независимой проверки показаны границами.}

Уточнённое решение совпадает с выбранной дугой Ламберта
в пределах численного сравнения. Это результат для заданных
начальных данных и времён. Уточнение конечного условия
не устанавливает минимальность характеристической скорости.

При добавлении возмущений меняются прогноз
$\vect y_f(\Delta\vect V)$ и его производные.
Манёвры конечной продолжительности требуют другой
параметризации управления и рассматриваются далее в курсе.

\clearpage
\begin{thebibliography}{9}
\addcontentsline{toc}{section}{Литература}
\bibitem{chiniaev}
Чиняев В.~Н., Кузнецов А.~А., Чугунов В.~В.
\textit{Ускоренный алгоритм решения задачи встречи
космических аппаратов на возмущённых эллиптических орбитах}.
Предоставленная автором редакция статьи.
Общая постановка: \S~2; линейная двухимпульсная задача:
\S~3; нелинейное уточнение: \S~4.2.
В этом конспекте общая схема записана в классических элементах
и при фиксированных временах манёвров.

\bibitem{bmw}
Bate R.~R., Mueller D.~D., White J.~E.
\textit{Fundamentals of Astrodynamics}.
New York: Dover, 1971.
Краевая задача по двум положениям и времени:
\S~5.2, с.~228--231; вариант итерации по фокальному
параметру: \S~5.4. Изменение орбитальных элементов:
\S~9.4.1, с.~401--407.
Фаза в настоящем тексте --- текущая средняя аномалия,
а не её значение на фиксированную эпоху.

\bibitem{izzo}
Izzo D. Revisiting Lambert's Problem //
\textit{Celestial Mechanics and Dynamical Astronomy}.
2015. Vol.~121. P.~1--15.
\href{https://doi.org/10.1007/s10569-014-9587-y}
{DOI: 10.1007/s10569-014-9587-y}.
Геометрия, классификация и многовитковые ветви
краевой задачи; вычислительная параметризация
этой работы здесь не используется.

\bibitem{cw}
Clohessy W.~H., Wiltshire R.~S.
Terminal Guidance System for Satellite Rendezvous //
\textit{Journal of the Aerospace Sciences}.
1960. Vol.~27, No.~9. P.~653--658.
\href{https://doi.org/10.2514/8.8704}{DOI: 10.2514/8.8704}.

\bibitem{ya}
Yamanaka K., Ankersen F.
New State Transition Matrix for Relative Motion
on an Arbitrary Elliptical Orbit //
\textit{Journal of Guidance, Control, and Dynamics}.
2002. Vol.~25, No.~1. P.~60--66.
\href{https://doi.org/10.2514/2.4875}{DOI: 10.2514/2.4875}.
Временные относительные уравнения: с.~61--62;
масштабирование и аналитическая матрица перехода:
с.~62--65. Конвенция осей первоисточника
отличается от принятой здесь RTN.

\bibitem{ss}
Schweighart S.~A., Sedwick R.~J.
High-Fidelity Linearized $J_2$ Model for Satellite Formation Flight //
\textit{Journal of Guidance, Control, and Dynamics}.
2002. Vol.~25, No.~6. P.~1073--1080.
\href{https://doi.org/10.2514/2.4986}{DOI: 10.2514/2.4986}.
Усреднённый градиент и опорная орбита: уравнения (13), (19);
относительная система: уравнение (45).

\bibitem{lange}
Lange B.~O., Smith R.~G.
\textit{The Application of Floquet Theory to the Computation
of Small Orbital Perturbations over Long Time Intervals
Using the Tschauner--Hempel Equations}.
Stanford University, SUDAER No.~241, 1965.
Нормированные уравнения TH: с.~3; вывод: приложение~A.
\href{https://ntrs.nasa.gov/citations/19660000761}{NASA NTRS}.
\end{thebibliography}

\end{document}
